%قبل از شروع مسئله ترجیحاً چیزی ننویسید. اما اگر خواستید بنویسید، در قالب لیستی (فقرات و شمارش) بنویسید.
%برای اضافه کردن نمره، بعدی از شروع‌مسئله، [x] را اضافه کنید و به جای x نمره مورد نظر را قرار دهید. برای زیرمسئله‌ها هم به همین صورت است.
%\شروع‌مسئله[10]
\شروع‌مسئله [20]
%\متن‌سیاه{سوال عملی اول:}\\
%\متن‌سیاه{دراپ‌تست}\\

% صورت اصلی مسئله را در این‌جا بنویسید.
شایان که به تازگی از جلسات مهمی اخراج شده است، در $n$ کلاس مشغول تدریس المپیاد است. در راستای جمع کردن کلاس‌ها با مساله‌ای روبه‌رو شده، از آنجایی که وقتی برای حل کردن سوال برایش باقی نمانده، از کیوان که در این روزها بسیار بیکار است می‌خواهد که مسأله زیر را برایش حل کند:\\*
در ابتدا کلاس‌های شایان خالی هستند، او $q$ درخواست دارد که به یکی از شکل‌های زیر می‌باشند.
\شروع{شمارش}
\فقره یک دانش‌آموز با نمره $x$ به کلاس $c$ اضافه شود. 
\فقره نمره تمام دانش‌آموزان کلاس $c$ با $x$ جمع شود.
\فقره دانش‌آموزی که کمترین نمره را در بین تمام دانش‌آموزان شایان دارد، اخراج کرده و نمره‌اش را گزارش کنید (اگر دانش‌آموزی وجود نداشت $-1$ را گزارش کنید).
\پایان{شمارش}
از آنجایی که کیوان دیگر توانایی‌های سابق را ندارد، از شما می‌خواهد تا به درخواست‌ها پاسخ دهید تا اندک آبرویی که برایش مانده حفظ شود.\\*
\متن‌سیاه{توجه کنید که نمرات شایان می‌تواند هر عدد صحیحی باشد (حتی منفی).}

\noindent
\متن‌سیاه{ورودی:}\\
%در خط اول ورودی، ابتدا k و سپس m داده می‌شود.\\
در خط اول ورودی $n$ و $q$ که تعداد کلاس‌ها و تعداد درخواست‌های شایان هستند، داده می‌شود.\\*
در $q$ خط بعدی به ترتیب درخواست‌ها به شکل $t_i$ و اگر $t_i \neq 3$ در ادامه دو عدد $c_i$ و $x_i$ نیز داده می‌شود که $t_i$ نوع درخواست $i$اُم می‌باشد.

\noindent
\متن‌سیاه{خروجی:}\\
%در تنها سطر خروجی، تعداد آیفون‌های مورد نیاز می‌آید. اگر در هیچ صورتی چنین کاری ممکن نبود ۱- چاپ می‌شود.\\
به ازای هر درخواست نوع $3$ عدد خواسته شده را در یک خط گزارش کنید.\\
\متن‌سیاه{محدودیت‌ها:}\\
$1\leq n,q\leq 10^5$\\
$1\leq t_i\leq 3$\\
$1\leq c_i\leq n$\\
$-10^9\leq x_i\leq 10^9$



%$m\leq10^9$

%\textbf{\متن‌سیاه{ورودی و خروجی نمونه}\\

\newpage
\متن‌سیاه{ورودی و خروجی نمونه:}\\


%% تست اول
\begin{latin}
	\begin{tabularx}{\textwidth}{|X | X|}
		\hline
		\textbf{stdin} &
		\textbf{stdout}\\
		\hline
		5 11
		
		1 1 5
		
		1 1 3
		
		1 2 2
		
		3
		
		3
		
		3
		
		3
		
		1 1 10
		
		1 2 5
		
		2 2 10
		
		3&
		2
		
		3
		
		5
		
		-1
		
		10\\
		\hline
		
	\end{tabularx}
\end{latin}

%% تست دوم

%اگر مسئله چند قسمت داشت به این صورت عمل کنید:
%\شروع{زیرمسائل}
%\زیرمسئله[0] % مسئلۀ ۱آ
%\زیرمسئله % مسئلۀ ۱ب
%\زیرمسئله % مسئلۀ ۱پ
%\زیرمسئله % مسئلۀ ۱ت
%\پایان{زیرمسائل}

\solution{
به ازای هر کلاس یک هرم در نظر می‌گیریم و یک هرم دیگر نیز داریم که کمترین نمره هر کلاس و به همراه شماره کلاس آن در این هرم قرار دارد. برای هر کلاس یک آرایه $toAdd$ نیز داریم که نشان می‌دهد تا الان چقدر درخواست نوع $2$ به نمرات این کلاس اضافه کرده. برای مدیریت درخواست نوع $1$، باید نمره $x-val[c]$ را به کلاس $c$ اضافه کنیم و برای مدیریت درخواست نوع $3$، باید همیشه هرمی که کمترین نمره کلاس‌ها را نگه می‌دارد، به‌روز نگه داریم، که این هم کار سختی نیست. در کل هر درخواست را توانستیم 
$\Theta(log q)$
پاسخ دهیم.\\
}




%%%%%%%%%%%%%%%%%%%%%%%%%%%%%%%%%%%%%%%%%%%%%%%%%%%%%
% برخی دستورات مفید برای تمرین‌ها. %
%%%%%%%%%%%%%%%%%%%%%%%%%%%%%%%%%%%%%%%%%%%%%%%%%%%%%

% فرمول‌های نمونه:
%$f(x) = \Theta(x)$
%$T(x, y) \leq \log(x) + 2^y + \binom{2n}{n}$
% {\tt code\_function}


% لیست موارد بدون شماره‌گذاری:
%\شروع{فقرات}
% \فقره مورد اول 
% \شروع{فقرات}
% \فقره مورد اول داخلی 
% \شروع{فقرات}
% \فقره اولین مورد خیلی داخلی 
% \فقره دومین مورد خیلی داخلی 
% \پایان{فقرات}
% \فقره مورد دوم داخلی 
% \پایان{فقرات}
% \فقره مورد دوم 
%\پایان{فقرات}

% لیست موارد شماره‌گذاری شده:
%\شروع{شمارش}
% \فقره یک مورد 
% \فقره یک مورد دیگر 
% \فقره باز هم یک مورد دیگر
%\پایان{شمارش}

% برای نوشتن یک فرمول چند خطی که یک شماره بخوره. برای حالت‌های دیگر به سایت زیر مراجعه کنید:
% https://www.overleaf.com/learn/latex/Aligning_equations_with_amsmath
%\begin{align} 
% \begin{split}
% (x+y)^3 &= (x+y)^2(x+y) \\
% &= (x^2+2xy+y^2)(x+y) \\
% &= (x^3+2x^2y+xy^2) + (x^2y+2xy^2+y^3) \\
% &= x^3+3x^2y+3xy^2+y^3
% \end{split} 
%\end{align}

% نوشتن ماتریس و آرایه دوبعدی
%\begin{align}
% A = 
% \begin{bmatrix}
% A_{11} & A_{21} \\
% A_{21} & A_{22}
% \end{bmatrix}
%\end{align}

% اضافه کردن تصویر با شرح. برای اضافه کردن تصویر بدون شرح، فقط دستور درج‌تصویر را استفاده کنید.
%\شروع{شکل}[h!]
%\تنظیم‌ازوسط
% \درج‌تصویر[پهنا=0.5\پهنای‌متن]{img.jpg}
%\شرح{شرح تصویر}
%\برچسب{تصویر:۱}
%\پایان{شکل}
%\پایان{زیرمسائل}

%کد زیر را ببینید.
%%\lstinputlisting{source_filename.py}
%\begin{latin}
%\begin{lstlisting}
%# Import the modules
%import sys
%import random
%
%ans = True
%
%while ans:
%question = raw_input("Ask the magic 8 ball a question: (press enter to quit) ")
%
%answers = random.randint(1,8)
%
%if question == "":
%sys.exit()
%
%elif answers == 1:
%print "It is certain"
%
%elif answers == 2:
%print "Outlook good"
%
%elif answers == 3:
%print "You may rely on it"
%
%elif answers == 4:
%print "Ask again later"
%
%elif answers == 5:
%print "Concentrate and ask again"
%
%elif answers == 6:
%print "Reply hazy, try again"
%
%elif answers == 7:
%print "My reply is no"
%
%elif answers == 8:
%print "My sources say no"
%\end{lstlisting}
%\end{latin}